% Source: https://artofproblemsolving.com/community/c3h2356858_why_does_this_happen
% Stable identity: aops-topic-2356858-q01
The error is the first sentence: the rule $a^{\dfrac bc}=\sqrt[c]{a^b}$ is only valid (as a definition of a real power that does not depend on how the exponent is written) for $a\ge0$, and it was used with the negative base $a=-1$.

For $a>0$, the rule gives the same value for every way of writing the exponent: $\sqrt[kc]{a^{kb}}$ and $\sqrt[c]{a^b}$ are both the positive number whose $kc$-th power is $a^{kb}$, so they are equal. For a negative base this fails, and the argument shows exactly that: the formula applied to $\dfrac35$ and to $\dfrac{6}{10}$ produces different numbers, so $(-1)^{\dfrac{6}{10}}=\sqrt[10]{(-1)^6}$ is not a legitimate step. (Squaring first, $(-1)^6=1$, destroys the sign; and $\left(\sqrt[10]{-1}\right)^6$ is not even a real number.)

The usual convention for a negative base is to allow a rational exponent only in lowest terms with an odd denominator, $a^{\dfrac pq}=\sqrt[q]{a^p}$ with $q$ odd. With it, $\dfrac{6}{10}$ must first be reduced to $\dfrac35$, and
\[(-1)^{\dfrac35}=(-1)^{\dfrac{6}{10}}=\sqrt[5]{(-1)^3}=\boxed{-1}.\]
So the step ``$=1$'' is the false one, and no contradiction arises.
